% ============================================================
% EDC EPISTEMIC MACROS (COMPLETE)
% ============================================================
% Version: 1.0 Final
% Date: 2026-01-30
% Purpose: LaTeX macros for consistent epistemic tagging
% Integration: Add to preamble of Part II main.tex
% ============================================================

% Required packages
\usepackage{xcolor}      % For colored tags (optional)
\usepackage{amsmath}     % For math environments
\usepackage{amssymb}     % For checkmark symbol

% ============================================================
% PRIMARY EPISTEMIC TAGS
% ============================================================

% Basic tags (no color)
\newcommand{\tagM}{\textbf{[M]}}
\newcommand{\tagP}{\textbf{[P]}}
\newcommand{\tagDer}{\textbf{[Der]}}
\newcommand{\tagDc}{\textbf{[Dc]}}
\newcommand{\tagI}{\textbf{[I]}}
\newcommand{\tagCal}{\textbf{[Cal]}}
\newcommand{\tagBL}{\textbf{[BL]}}
\newcommand{\tagOpen}{\textbf{[Open]}}

% Colored tags (optional - uncomment if desired)
% \renewcommand{\tagM}{\textcolor{blue}{\textbf{[M]}}}
% \renewcommand{\tagP}{\textcolor{purple}{\textbf{[P]}}}
% \renewcommand{\tagDer}{\textcolor{green!60!black}{\textbf{[Der]}}}
% \renewcommand{\tagDc}{\textcolor{orange!80!black}{\textbf{[Dc]}}}
% \renewcommand{\tagI}{\textcolor{cyan!60!black}{\textbf{[I]}}}
% \renewcommand{\tagCal}{\textcolor{red!70!black}{\textbf{[Cal]}}}
% \renewcommand{\tagBL}{\textcolor{gray}{\textbf{[BL]}}}
% \renewcommand{\tagOpen}{\textcolor{red}{\textbf{[Open]}}}

% ============================================================
% SUBTAGS (PRECISION REFINEMENT)
% ============================================================

% Derivation subtags
\newcommand{\tagDerSym}{\textbf{[Der:Sym]}}
\newcommand{\tagDerNum}{\textbf{[Der:Num]}}

% Conditional derivation subtags
\newcommand{\tagDcSym}{\textbf{[Dc:Sym]}}
\newcommand{\tagDcNum}{\textbf{[Dc:Num]}}
\newcommand{\tagDcApprox}{\textbf{[Dc:Approx]}}

% ============================================================
% REDUCTION NORMALIZATION FACTOR
% ============================================================

% Generic reduction factor
\newcommand{\Cred}{C_{\mathrm{red}}}

% Specific reduction factors (by quantity)
\newcommand{\Credm}{C_{\mathrm{red}}^{(m)}}           % mass ratio
\newcommand{\Credalpha}{C_{\mathrm{red}}^{(\alpha)}}  % fine structure
\newcommand{\CredGF}{C_{\mathrm{red}}^{(G_F)}}        % Fermi coupling
\newcommand{\CredW}{C_{\mathrm{red}}^{(\theta_W)}}    % Weinberg angle

% Generic parameterized version
\newcommand{\Credgen}[1]{C_{\mathrm{red}}^{(#1)}}     % C_red^{(label)}

% ============================================================
% REFERENCE MACROS
% ============================================================

% Reference to baseline constants table
\newcommand{\refBaseline}{Table~\ref{tab:baseline_constants}}

% Reference to epistemic standard section
\newcommand{\refEpistemic}{Section~\ref{sec:epistemic_standard}}

% Reference to open problems register
\newcommand{\refOPR}{Appendix~\ref{app:opr}}

% Generic OPR reference
\newcommand{\OPR}[1]{OPR-#1}  % Usage: \OPR{21} → OPR-21

% ============================================================
% STATUS BOXES
% ============================================================

% Requires tcolorbox package:
% \usepackage{tcolorbox}

% Status box environment
\newenvironment{statusbox}[1]{%
  \begin{tcolorbox}[colback=blue!5,colframe=blue!75!black,title=#1]
}{%
  \end{tcolorbox}
}

% Warning box environment
\newenvironment{warningbox}[1]{%
  \begin{tcolorbox}[colback=red!5,colframe=red!75!black,title=#1]
}{%
  \end{tcolorbox}
}

% Open problem box environment
\newenvironment{openbox}[1]{%
  \begin{tcolorbox}[colback=orange!5,colframe=orange!75!black,title={Open: #1}]
}{%
  \end{tcolorbox}
}

% ============================================================
% COMMON ABBREVIATIONS
% ============================================================

% Observables
\newcommand{\mpme}{\frac{m_p}{m_e}}           % Proton-electron mass ratio
\newcommand{\alphaInv}{\alpha^{-1}}            % Fine structure inverse
\newcommand{\sinThetaW}{\sin^2\theta_W}        % Weinberg angle
\newcommand{\DeltaMnp}{\Delta m_{np}}          % Neutron-proton mass diff

% Common symbols
\newcommand{\Lzero}{L_0}                       % Junction extent
\newcommand{\ellp}{\ell_p}                     % Proton length
\newcommand{\Rxi}{R_\xi}                       % Compactification radius
\newcommand{\tauN}{\tau_n}                     % Neutron lifetime

% ============================================================
% COMPARISON MACROS
% ============================================================

% Comparison statement (EDC vs BL)
\newcommand{\comparison}[4]{%
  % Usage: \comparison{quantity}{EDC_value}{BL_value}{error_percent}
  \begin{align*}
  \text{EDC:} \quad & #1 = #2 \\
  \text{[BL]:} \quad & #1 = #3 \\
  \text{Error:} \quad & #4\%
  \end{align*}
}

% ============================================================
% EPISTEMIC STATEMENT MACROS
% ============================================================

% Conditional statement (for [Dc] results)
\newcommand{\conditional}[1]{%
  \textit{Conditional on: #1}
}

% Mathematical precision statement
\newcommand{\precision}[1]{%
  \textit{Precision: #1}
}

% ============================================================
% EXAMPLE USAGE
% ============================================================

% Example 1: Basic tagging
% Result: $m_p/m_e = 6\pi^5$ \tagDcSym

% Example 2: Reduction factor
% Normalization: $\Credm \approx 1.0002$

% Example 3: Comparison
% \comparison{m_p/m_e}{6\pi^5 = 1836.118}{1836.153}{0.002}

% Example 4: Reference
% See \refBaseline for baseline data.

% Example 5: OPR reference
% This depends on \OPR{21} (brane boundary value problem).

% Example 6: Status box
% \begin{statusbox}{Current Status}
% Result is \tagDcSym with 0.002\% agreement.
% \end{statusbox}

% ============================================================
% End of Macros
% ============================================================
